\documentclass[12pt,a4paper]{article}
%documentclass[asa,11pt]{article}
\usepackage[utf8]{inputenc}
\usepackage[margin=1.2in]{geometry} 

\usepackage{amsmath}
\usepackage{graphicx}
\usepackage{amsmath}
\usepackage{tipa, amssymb}
\usepackage{float}
\usepackage{accents}
%\newcommand{\undertilde}[1]{\underset{\widetilde{}}{#1}}
\usepackage{graphicx} % leave off '[demo]' in real document\graphicspath{ {./images/} }
%%%%%%%%%%%%%%%%%%%%%%%%%%%%%%%%%%%%%%%%%%%%%%%%%%%%%%%%%%%%%%%%%%%%%%%%%%
% Balgobin Nandram
\newbox\TempBox \newbox\TempBoxA
\def\uwd#1{%
	\ifmmode\setbox\TempBox=\hbox{$#1$}\else\setbox\TempBox=\hbox{#1}\fi%
	\setbox\TempBoxA=\hbox to \wd\TempBox{\hss\char'176\hss}%
	\rlap{\copy\TempBox}\smash{\lower10pt\hbox{\copy\TempBoxA}}%
}
\def\mathunderaccent#1{\let\theaccent#1\mathpalette\putaccentunder}
\def\putaccentunder#1#2{\oalign{$#1#2$\crcr\hidewidth
		\vbox to.2ex{\hbox{$#1\theaccent{}$}\vss}\hidewidth}}
\def\ttilde#1{\tilde{\tilde{#1}}}
%%%%%%%%%%%%%%%%%%%%%%%%%%%%%%%%%%%%%%%%%%%%%%%%%%%%%%%%%%%%%%%%%%%%%%%%%%%%
\title{Supplemental Material}
\begin{document}

\begin{center}
{\Large {\bf Supplemental Materials (SM), July 5, 2026}}
\end{center}

%\begin{center}
\section*{Appendix C: Identifiability in the Mixture Model}
%\end{center}
%\def\theequation{G.\arabic{equation}}
%\setcounter{equation}{0}

We have $(\uwd{x}_i, y_i), i=1,\ldots,n,$  where
$\uwd{x}_i$ are covariates and $y_i$ are values of the study variable, 
and we consider the mixture model,
$$
y_i \mid p,\uwd{\beta}, \sigma^2, \rho \stackrel{ind} \sim (1-p)\mbox{Normal}\left(\uwd{x}_i^\prime 
\uwd{\beta}, \rho \sigma^2\right) + p \mbox{Normal}\left(\uwd{x}_i^\prime \uwd{\beta}, \sigma^2\right), 0< p < \frac{1}{2}, 0<\rho < 1.
$$
Without loss of generality, we drop the survey weights.
Note that for each $i=1,\ldots,n,$ the two components have the same mean, $\uwd{x}_i^\prime \uwd{\beta}$, so that the distribution is symmetric about the mean, and outliers can be in both
tails, although we can consider outliers only in the right tail. Here, we show that the
parameters $p, \uwd{\beta}, \rho, \sigma^2$ are identifiable. We also show that
the joint posterior density of $\pi(p, \uwd{\beta},\rho,\sigma^2\mid \uwd{y})$
under flat priors is proper.

Identifiability is demonstrated in the following steps.
\begin{itemize}
   \item[1.] Separate the data into two groups, $G$ and $G^c$. This can be done by
	ordering the $y_1,\ldots,y_n$ from smallest to largest and looking for a separation
	among the larger values. Let $m > \frac{n}{2}$ denote the cardinality of $G$. Then
	$\hat{p} = (n-m)/n$.
	 \item[2.] Since the two components of the $y_i$ have the same mean, we can use least squares
	 to get $\hat{\uwd{\beta}} = \biggl(\sum_{i=1}^n \uwd{x}_i \uwd{x_i}^\prime\biggr)^{-1} 
	 \sum_{i=1}^n \uwd{x}_i y_i$, where we assume the design matrix is full rank.
	 \item[3.] Consider the residuals, $y_i -\uwd{x}_i^\prime \hat{\uwd{\beta}}, i=1,\ldots,n$. 
	  Using the residuals in $G$, we can obtain the intra-class correlation, $\hat{\rho}$. This can be obtained by pairing the residuals. If $\hat{\rho} < 0$, take $\hat{\rho} = 0$. Also, let $S_0^2$ denote the sample variance of the
			residuals in $G$ and $S_1^2$ the sample variance of the residuals in $G^c$. Then,
			$\hat{\rho} = \frac{S_0^2}{S_1^2}$, and 
			 if $\hat{\rho} > 1$, take $\hat{\rho} = 1$.
	 \item[4.] Noting that $\mbox{Var}(y_i) = \left\{(1-p)\rho + p \right\} \sigma^2, i=1,\ldots,n$, we have
	$$
	\left\{(1-\hat{p}) \hat{\rho} +\hat{p} \right\} \hat{\sigma^2} = \frac{\sum_{i=1}^n(y_i-\uwd{x}_i^\prime \uwd{\beta})^2}{n-c},
	$$
	where $c$ is the number covariates including the intercept. That is,
	$$
	\hat{\sigma^2} = \frac{\sum_{i=1}^n(y_i-\uwd{x}_i^\prime \uwd{\beta})^2}{(n-c)	
	\{(1-\hat{p}) \hat{\rho} +\hat{p} \}}. 
	$$
\end{itemize} 

Therefore, the parameters, $p, \uwd{\beta}, \rho, \sigma^2$, are identifiable. Of course, fitting the Bayesian model with prior,
$$
\pi(p,\uwd{\beta},\rho, \sigma^2) \propto \frac{1}{\sigma^2},
0<p<\frac{1}{2}, 0 <\rho< 1,
$$
will provide posterior densities of these parameters, and better posterior summaries.

%\newpage

%\begin{center}
\section*{Appendix D:  Propriety of the Joint Posterior Density of the Mixture Model}
%\end{center}
%\def\theequation{H.\arabic{equation}}
%\setcounter{equation}{0}

We consider the Bayesian version of our mixture model,
$$
y_i \mid p,\uwd{\beta}, \sigma^2, \rho \stackrel{ind} \sim (1-p)\mbox{Normal}\left(\uwd{x}_i^\prime 
\uwd{\beta}, \rho \sigma^2\right) + p \mbox{Normal}\left(\uwd{x}_i^\prime \uwd{\beta}, \sigma^2\right),
i=1,\ldots,n, 
$$
$$
\pi(p,\uwd{\beta},\rho,\sigma^2) \propto \frac{1}{\sigma^2}, 0<p<\frac{1}{2},
0 < \rho < 1.
$$
Without loss of generality, we drop the survey weights.
The joint posterior density is
$$
\pi(p, \uwd{\beta},\rho,\sigma^2 \mid \uwd{y}) \propto
$$
$$
\frac{1}{\sigma^2} \prod_{i=1}^n\biggl\{(1-p)\mbox{Normal}\left(\uwd{x}_i^\prime \uwd{\beta},
\rho \sigma^2\right) + p\mbox{Normal}\left(\uwd{x}_i^\prime \uwd{\beta},
\sigma^2\right)\biggr\}, 0<p<\frac{1}{2}, 0<\rho<1.
$$
It is easy to show that the joint posterior density is proper. 

For this purpose, we introduce independent binary latent variables, $z_i = 0, 1, i=1,\ldots,n$. Then, $(y_i,z_i), i=1,\ldots,n,$ are independent, and the joint posterior density is
$$
\pi(\uwd{z}, p, \uwd{\beta},\rho,\sigma^2 \mid \uwd{y}) \propto
$$
$$
\frac{1}{\sigma^2} 
\times
\prod_{i=1}^n \biggl\{(1-p)\mbox{Normal}\left(\uwd{x}_i^\prime \uwd{\beta},
\rho \sigma^2\right)\biggr\}^{1-z_i} \biggl\{p\mbox{Normal}\left(\uwd{x}_i^\prime \uwd{\beta},
\sigma^2\right)\biggr\}^{z_i}, 0<p<\frac{1}{2}, 0<\rho<1.
$$
Assuming that the design matrix is full rank,
it is easy to integrate out $(p, \uwd{\beta}, \sigma^2)$ to get $\pi(\uwd{z},\rho\mid \uwd{y})$.

We can now sum $\pi(\uwd{z},\rho\mid \uwd{y})$ over all $z_i$ ($n$ summations) to get
$\pi(\rho \mid \uwd{y})$.
Then, finally assuming $\epsilon \leq \rho \leq 1-\epsilon$, where $\epsilon$ is
a small positive real number, the 
integral over $\rho$ is finite. Therefore,
$\pi(p, \uwd{\beta},\rho,\sigma^2\mid \uwd{y})$ is proper.

\clearpage
\section*{Appendix E: Summary results of population parameters}
%\clearpage

\setcounter{table}{2}
\begin{center}
\begin{minipage}{0.8\textwidth}
    \begin{table}[H]
        \centering
        \caption{Summaries of the parameters of non-robust population model.}
        \resizebox{\textwidth}{!}{%
            \begin{tabular}{@{}l*{6}{c}@{}} 
                \hline
                & Summaries & $\beta_0$ & $\beta_1$ & $\beta_2$ & $\beta_3$ & $\sigma^2$ \\ 
                \hline
                TR & ESS & 1000 & 1000 & 1000 & 1000 & 1000 \\
                   & p-value & 0.140 & 0.297 & 0.633 & 0.053 & 0.521 \\
                UNP & ESS & 1000 & 1000 & 1000 & 1000 & 1000 \\
                    & p-value & 0.649 & 0.566 & 0.667 & 0.681 & 0.556 \\
                UNNP & ESS & 1000 & 1000 & 1000 & 1000 & 1000 \\
                     & p-value & 0.755 & 0.605 & 0.933 & 0.639 & 0.713 \\
                \hline
            \end{tabular}%
        }
        \vspace{0.5em}
        \begin{minipage}{0.95\textwidth}
            \footnotesize
            \textit{Note:} Summaries are Effective Sample Size (ESS) and p-value of Geweke test.
        \end{minipage}
    \end{table}
\end{minipage}

\vspace{3cm}

\begin{minipage}{\textwidth}
    \begin{table}[H]
        \centering
        \caption{Posterior summaries of non-robust population model.}
        \resizebox{\textwidth}{!}{%
            \begin{tabular}{@{}l*{6}{c}@{}} 
                \hline
                & Summaries & $\beta_0$ & $\beta_1$ & $\beta_2$ & $\beta_3$ & $\sigma^2$ \\
                \hline
                TR & PM & 27.613 & -0.012 & 0.647 & 0.482 & 15.079 \\
                   & PSD & 0.451 & 0.007 & 0.453 & 0.261 & 0.549 \\
                   & NSE & 0.0002 & 0.0142 & 0.008 & 0.014 & 0.017 \\
                   & PCV & 0.016 & 0.583 & 0.7 & 0.541 & 0.036 \\
                   & $95\%$ HPDI & (26.773,28.507) & (-0.025,0.001) & (-0.230,1.529) & (0.003,1.011) & (14.06,16.14) \\
                UNP & PM & 28.909 & -0.029 & 0.494 & 0.429 & 15.278 \\
                    & PSD & 2.609 & 0.039 & 0.874 & 0.345 & 7.598 \\
                    & NSE & 0.001 & 0.020 & 0.010 & 0.066 & 0.221 \\
                    & PCV & 0.090 & 1.344 & 1.769 & 0.804 & 0.497 \\
                    & $95\%$ HPDI & (25.725,35.326) & (-0.120,0.026) & (-0.927,1.661) & (-0.277,0.990) & (9.277,22.660) \\
                UNNP & PM & 28.473 & -0.022 & 0.498 & 0.467 & 16.281 \\
                     & PSD & 2.098 & 0.032 & 0.633 & 0.326 & 6.991 \\
                     & NSE & 0.001 & 0.020 & 0.010 & 0.066 & 0.221 \\
                     & PCV & 0.074 & 1.455 & 1.271 & 0.698 & 0.429 \\
                     & $95\%$ HPDI & (25.659,33.670) & (-0.099,0.250) & (-0.743,1.561) & (-0.191,1.072) & (10.689,22.137) \\
                \hline
            \end{tabular}%
        }
        \vspace{0.5em}
        \begin{minipage}{0.95\textwidth}
            \footnotesize
            \textit{Note:} Summaries are PM = posterior mean, PSD = posterior standard deviation, NSE = numeric standard error, PCV = posterior coefficient of variation, HPDI = highest posterior density interval.
        \end{minipage}
    \end{table}
\end{minipage}
\end{center}

%\clearpage

% ---------- PAGE 2 ----------
\begin{center}
\begin{minipage}{\textwidth}
    \begin{table}[H]
        \centering
        \caption{Summaries of the parameters of robust population model.}
        \resizebox{\textwidth}{!}{%
            \begin{tabular}{@{}l*{8}{c}@{}} 
                \hline
                & Summaries & $\beta_0$ & $\beta_1$ & $\beta_2$ & $\beta_3$ & $\sigma^2$ & $\rho$ & $p$ \\
                \hline
                TR-R & ESS & 1000 & 1000 & 1000 & 1000 & 1000 & 1000 & 1000 \\
                      & p-value & 0.902 & 0.862 & 0.218 & 0.557 & 0.413 & 0.957 & 0.053 \\
                UNP-R & ESS & 1000 & 1000 & 1000 & 1000 & 1000 & 1000 & 1000 \\
                       & p-value & 0.487 & 0.575 & 0.402 & 0.312 & 0.858 & 0.856 & 0.684 \\
                UNNP-R & ESS & 1000 & 1000 & 1000 & 1000 & 1000 & 1000 & 1000 \\
                        & p-value & 0.736 & 0.825 & 0.556 & 0.113 & 0.737 & 0.822 & 0.157 \\
                \hline
            \end{tabular}%
        }
        \vspace{0.5em}
        \begin{minipage}{0.95\textwidth}
            \footnotesize
            \textit{Note:} Summaries are Effective Sample Size (ESS) and p-value of Geweke test.
        \end{minipage}
    \end{table}
\end{minipage}

\vspace{3cm}

\begin{minipage}{1.1\textwidth}
    \begin{table}[H]
        \centering
        \caption{Posterior summaries of robust population model.}
        \resizebox{\textwidth}{!}{%
            \begin{tabular}{@{}l*{8}{c}@{}} 
                \hline
                & Summaries & $\beta_0$ & $\beta_1$ & $\beta_2$ & $\beta_3$ & $\sigma^2$ & $\rho$ & $p$ \\
                \hline
                TR-R & PM & 26.940 & -0.001 & 0.030 & -0.024 & 53.19 & 0.167 & 0.159 \\
                     & PSD & 0.433 & 0.007 & 0.400 & 0.260 & 12.334 & 0.033 & 0.040 \\
                     & NSE & 0.013 & 0.0002 & 0.012 & 0.008 & 0.391 & 0.001 & 1.971 \\
                     & PCV & 0.016 & 7.00 & 13.3 & 10.8 & 0.231 & 0.197 & 0.251 \\
                     & $95\%$ HPDI & (26.095,27.809) & (-0.015,0.012) & (-0.773,0.790) & (-0.515,0.495) & (37.191,83.829) & (0.102,0.236) & (0.084,0.236) \\
                UNP-R & PM & 27.425 & -0.008 & 0.171 & 0.034 & 35.972 & 0.137 & 0.313 \\
                       & PSD & 1.707 & 0.025 & 0.688 & 0.314 & 16.609 & 0.048 & 0.155 \\
                       & NSE & 0.053 & 0.0007 & 0.021 & 0.009 & 0.513 & 0.001 & 0.004 \\
                       & PCV & 0.062 & 3.125 & 5.70 & 9.235 & 0.462 & 0.35 & 0.495 \\
                       & $95\%$ HPDI & (25.284,31.656) & (-0.070,0.024) & (-0.932,1.366) & (-0.574,0.631) & (6.392,68.697) & (0.038,0.222) & (0.031,0.499) \\
                UNNP-R & PM & 27.209 & -0.005 & 0.228 & 0.031 & 37.399 & 0.148 & 0.305 \\
                        & PSD & 1.564 & 0.023 & 0.609 & 0.292 & 15.662 & 0.048 & 0.140 \\
                        & NSE & 0.048 & 0.0007 & 0.019 & 0.009 & 0.498 & 0.001 & 0.004 \\
                        & PCV & 0.057 & 4.60 & 0.57 & 0.119 & 0.418 & 0.324 & 0.459 \\
                        & $95\%$ HPDI & (25.239,30.635) & (-0.057,0.027) & (-0.866,1.377) & (-0.551,0.636) & (9.686,68.290) & (0.048,0.237) & (0.051,0.498) \\
                \hline
            \end{tabular}%
        }
        \vspace{0.5em}
        \begin{minipage}{0.95\textwidth}
            \footnotesize
            \textit{Note:} Summaries are PM = posterior mean, PSD = posterior standard deviation, NSE = numeric standard error, PCV = posterior coefficient of variation, HPDI = highest posterior density interval.
        \end{minipage}
    \end{table}
\end{minipage}
\end{center}

%%%%%%%%%%%%%%%%%%%%%%%%%%%%%%%%%%%%%%%%%%%%%%%%%%%%%%%%%%%%%%%%%%%%%%%%%%%%%%%%%%%%

\newpage

\section*{Appendix F: Boxplots for adjusted weights of all units in non-probability sample}


\setcounter{figure}{3}

In Figures 1 and 2, we show posterior densities of $\Bar{Y}$ under the six models.
\begin{figure}[ht]
	\centering
	\includegraphics[width=14.8cm]{y_bar_compare.png}
	\caption{Posterior distributions of $\Bar{Y}$ for TR(1), UNP(2) and UNNP(3).}
	\label{fig1}
\end{figure}
\begin{figure}[h]
	\centering
	\includegraphics[width=14.8cm]{y_bar_compare_robust.png}
	\caption{Posterior distributions of $\Bar{Y}$ for TR-R(4), UNP-R(5) and UNNP-R(6).}
	\label{fig1}
\end{figure}
%%%%%%%%%%%%%%%%%%%%%%%%%%%%%%%%%%%%%%%%%%%%%%%%%%%%%%%%%%%%%%%%%%%%%%%%%%%%%%%%%%%


\clearpage



%
%\begin{figure}[ht]
	%\centering
	%\includegraphics[width=14.8cm]{y_bar_compare.png}
	%\caption{Posterior distributions of $\Bar{Y}$ for TR(1), UNP(2) and UNNP(3).}
	%\label{fig1}
%\end{figure}
%\begin{figure}[h]
	%\centering
	%\includegraphics[width=14.8cm]{y_bar_compare_robust.png}
	%\caption{Posterior distributions of $\Bar{Y}$ for TR-R(4), UNP-R(5) and UNNP-R(6).}
	%\label{fig1}
%\end{figure}
%












\section*{Appendix G: Boxplots for adjusted weights of all units in non-probability sample}
In Appendix G, we randomly select 500 units from the remaining observations, excluding the first 100 units in the non-probability sample. The boxplots represent
the posterior density of each adjusted weight, obtained from the 1,000 iterates
of the Gibbs sampler. We can see that some adjusted weights are very small while others are very large. Also, some
adjusted weights are much more variable than others.

%\setcounter{figure}{4}

\begin{figure}[h]
	\centering
	\includegraphics[width=14.8cm]{boxplot_part_5.png}
	\caption{Boxplots for units in non-probability sample.}
	\label{fig1}
\end{figure}

\begin{figure}[h]
	\centering
	\includegraphics[width=14.8cm]{boxplot_part_7.png}
	\caption{Boxplots for units in non-probability sample.}
	\label{fig1}
\end{figure}


\begin{figure}[h]
	\centering
	\includegraphics[width=14.8cm]{boxplot_part_10.png}
	\caption{Boxplots for units in non-probability sample.}
	\label{fig1}
\end{figure}

\begin{figure}[h]
	\centering
	\includegraphics[width=14.8cm]{boxplot_part_12.png}
	\caption{Boxplots for units in non-probability sample.}
	\label{fig1}
\end{figure}


\begin{figure}[h]
	\centering
	\includegraphics[width=14.8cm]{boxplot_part_15.png}
	\caption{Boxplots for units in non-probability sample.}
	\label{fig1}
\end{figure}
\end{document}